\documentclass[10pt,a4paper]{amsart}
\usepackage[margin=19mm]{geometry}
\usepackage{amsmath,amssymb,mathtools}
\usepackage[scr=rsfs,cal=euler]{mathalfa}
\usepackage{xfrac}
\usepackage[unicode,hidelinks,bookmarks=false]{hyperref}
\newcommand{\ie}{i.e.\ }
\newcommand{\cf}{cf.\ }
\newcommand{\resp}{respectively\ }
\newcommand{\Subsection}{\subsection}
\providecommand{\nobreakdash}{\nobreak\hbox{-}}
\setlength{\emergencystretch}{2em}
\multlinegap0pt
\usepackage{amssymb}
\newtheorem{lemma}{Lemma}
\newtheorem{theorem}{Theorem}
\newcommand{\OF}{\mathcal{O}_F}
\newcommand{\OK}{\mathcal{O}_K}
\newcommand{\Q}{\mathbb{Q}}
\newcommand{\Z}{\mathbb{Z}}
\title{\textbf{On the Natural Density of Monic Integer Polynomials \\ with Roots in a Fixed Number Field}}
\author{Amirali Fatehizadeh}
\date{Source: arXiv:2605.21513v1 (CC BY 4.0)}
\begin{document}
\maketitle

\noindent\emph{Source and licence.} \textbf{On the Natural Density of Monic Integer Polynomials \\ with Roots in a Fixed Number Field}. Version arXiv:2605.21513v1 (CC BY 4.0), distributed by arXiv under the Creative Commons Attribution 4.0 International licence, http://creativecommons.org/licenses/by/4.0/, as recorded in the arXiv licence metadata for this exact version and shown on the abstract page https://arxiv.org/abs/2605.21513v1. Full licence notice: \texttt{licenses/arxiv-2605.21513-CC-BY-4.0.txt}.\par\smallskip

\noindent\textit{Editorial scope.} Counted unit: the article's theorem thm:main (source0.tex lines 182--184). The article's complete original proof of it is delivered with it (source0.tex lines 186--247, one \texttt{proof} environment of 62 lines). No mathematics is written, repaired or re-derived: the statement and the proof are the article's own lines, in the article's own notation, and no retained line is substituted or re-flowed.  Retained uncounted as the source-local context the proof needs: lines 123--129.  Nothing was omitted from any retained span, so the package carries no omission record.  No retained line contains a citation, so the wrapper carries no bibliography and no bibliography entry.  No retained line contains a figure environment, an image-inclusion command or drawing code, so no image is delivered and the package declares no asset dependency.  Editorial formatting only, in addition: an AMS \texttt{amsart} preamble that restates the article's own declarations and macros used by the retained lines, namely the environments \texttt{lemma}, \texttt{theorem} on separate counters, together with the standard packages those lines need.  The retained environments print in this document as Lemma 1, Theorem 1 in order of appearance, whereas the article prints the same items with the numbering of its own full document.  The complete native source of the article as distributed in the arXiv e-print for this exact version is bundled unchanged as \texttt{sources/201066/source0.tex}.
\begin{lemma}[Landau-Mignotte Bound]\label{lem:landau}
Let $P \in \Z[x]$ be a polynomial of degree $n \ge 1$ and let $P(x) = m(x) \cdot Q(x)$ be a factorization, where $m, Q \in \Z[x]$ are polynomials with non-zero leading coefficients. Then there exists a constant $C_n \ge 1$ (depending only on $n$) such that
\begin{equation*}
    H(m) \cdot H(Q) \le C_n H(P).
\end{equation*}
Throughout this paper, we adopt the explicit value $C_n = 2^n \sqrt{n+1}$ as a standard bound.
\end{lemma}

\begin{theorem}\label{thm:main}
Let $K$ be a fixed number field. The natural density of the set of monic integer polynomials of degree $n$ that have at least one root in $K$ is zero.
\end{theorem}

\begin{proof}
Let $M_K(H) \subseteq \mathcal{B}_n(H)$ denote the set of all target polynomials. As discussed in the previous section, the existence of a root $\alpha \in K$ for a polynomial $P \in M_K(H)$ implies that $\alpha \in \OK$, and its minimal polynomial, $m_\alpha(x)$, must be a factor of $P(x)$ in the ring $\Z[x]$.

We partition the set $M_K(H)$ into two disjoint subsets based on its reducibility behavior over the field of rational numbers:
\begin{equation*}
    M_K(H) = \mathrm{red}_n(H) \cup \mathrm{irr}_n(H),
\end{equation*}
where $\mathrm{red}_n(H)$ comprises the reducible polynomials and $\mathrm{irr}_n(H)$ comprises the irreducible ones. We will bound the cardinality of each of these sets independently.

First, consider a polynomial $P \in \mathrm{red}_n(H)$. Such a polynomial admits a factorization $P(x) = m(x) \cdot Q(x)$ over $\Z[x]$, where $m(x)$ is the minimal polynomial of a root $\alpha \in \OK$ with degree $\deg(m) = d$. Since $P$ is reducible, the degree of this factor must lie in the range $1 \le d \le \min(n-1, D)$, where $D =[K:\Q]$. 

By Lemma \ref{lem:landau}, the height of the quotient $Q(x)$, which is a monic polynomial of degree $n-d$, satisfies the following inequality:
\begin{equation*}
    H(Q) \le \frac{C_n H(P)}{H(m)} \le \frac{C_n H}{H(m)}.
\end{equation*}
To precisely count the possible pairs $(m, Q)$, assume that the height of the minimal polynomial is fixed at $H(m) = k$.
\begin{itemize}
    \item \textbf{Number of choices for $m(x)$:} The monic polynomial $m(x)$ is determined by its $d$ non-leading coefficients. For its height to be exactly $k$, all its coefficients must lie in the interval $[-k, k]$, with at least one coefficient being exactly $\pm k$. The total number of such polynomials is exactly $(2k+1)^d - (2k-1)^d$, which exhibits an asymptotic behavior of $O(k^{d-1})$.
    \item \textbf{Number of choices for $Q(x)$:} Similarly, the monic polynomial $Q(x)$ is determined by its $n-d$ non-leading coefficients, each of which is bounded within the interval $[-C_n H/k, C_n H/k]$. Thus, the number of possible choices for $Q(x)$ is of the order $O\left( \left(\frac{H}{k}\right)^{n-d} \right)$.
\end{itemize}

By summing the product of these bounds over all possible heights $1 \le k \le C_n H$ and all permissible degrees $1 \le d \le \min(n-1, D)$, we obtain the following upper bound for the cardinality of the reducible component:
\begin{equation*}
    \#\mathrm{red}_n(H) \ll \sum_{d=1}^{\min(n-1, D)} \sum_{k=1}^{C_n H} k^{d-1} \cdot \left(\frac{H}{k}\right)^{n-d} = \sum_{d=1}^{\min(n-1, D)} H^{n-d} \sum_{k=1}^{C_n H} k^{2d-n-1}.
\end{equation*}

To evaluate the asymptotic behavior of this expression, we consider the exponent of the variable $k$ in the inner sum, denoted by $p = 2d - n - 1$. Note that $p$ is an integer. The asymptotic behavior of the series $\sum k^p$ depends entirely on the sign and value of $p$. We analyze this series in three distinct cases:

\begin{itemize}
    \item \textbf{Case 1 ($p < -1$):} In this regime, the series $\sum k^p$ converges. Consequently, the sum remains bounded by a constant independent of $H$, i.e., $O(1)$, regardless of how large $H$ becomes. Thus, the total contribution of this part is $H^{n-d} \cdot O(1) \ll H^{n-1}$.
    
    \item \textbf{Case 2 ($p = -1$):} This case occurs only if $n$ is even and $d = n/2$. Here, the sum reduces to a harmonic series, which exhibits logarithmic divergence, namely $\sum_{k=1}^{C_n H} k^{-1} = O(\log H)$. The contribution of this section is $H^{n-d}\log H$. For the specific case $n=2$ (where $d=1$), this yields a contribution of $O(H \log H)$. For $n \ge 4$, since $d = n/2 \le n-2$, this contribution is also dominated by $O(H^{n-1})$.
    
    \item \textbf{Case 3 ($p > -1$):} Under this condition, the series diverges, and its asymptotic behavior, based on an integral approximation, is $O(H^{p+1}) = O(H^{2d-n})$. Therefore, the contribution of this part is $H^{n-d} \cdot H^{2d-n} = H^d$. Since the maximum possible value for $d$ is $n-1$, this case is also strictly bounded by $O(H^{n-1})$.
\end{itemize}

Combining these three cases, we conclude that:
\begin{equation*}
    \#\mathrm{red}_n(H) \ll 
    \begin{cases} 
      H \log H & \text{if } n = 2, \\
      H^{n-1} & \text{if } n > 2.
    \end{cases}
\end{equation*}

Now, consider the case where $P \in \mathrm{irr}_n(H)$. In this scenario, the polynomial $P$ is irreducible over $\Q$ and possesses a root $\alpha \in \OK$. Consequently, $P$ must be the minimal polynomial of $\alpha$, and its degree is exactly $n$. The root $\alpha$ generates a subfield $F = \Q(\alpha) \subseteq K$ with degree $[F:\Q] = n$. Since the base field $K$ is fixed, the number of such subfields is finite; we denote the set of these subfields by $\mathrm{Sub}_n(K)$. Note that if $[K:\Q] < n$, the set of irreducible polynomials satisfying this condition is empty. Furthermore, each minimal polynomial of degree $n$ produces exactly $n$ conjugate roots. Therefore, counting the number of permissible elements $\alpha$ provides a valid upper bound for the number of these polynomials.

We first perform the count for a fixed subfield $F \in \mathrm{Sub}_n(K)$. The element $\alpha$ generates a principal ideal $\mathfrak{a} = (\alpha)$ in the ring of integers $\OF$. The algebraic norm of this root is equal to the absolute value of the constant term of the polynomial, i.e., $|N_{F/\Q}(\alpha)| = |P(0)| \le H$. As established in the previous section, the total number of possible principal ideals satisfying this norm constraint is at most of order $O(H)$.

Any other generator for these ideals is of the form $\alpha' = \alpha u$, where $u \in \OF^\times$ is a unit. By Dirichlet's Unit Theorem, since $[F:\Q] = n$, the rank of the unit group is $r = r_1 + r_2 - 1 \le n - 1$. By applying the logarithmic embedding, and given that all coefficients of $P$ are bounded by $H$, the logarithmic vector of the roots is confined within a bounded region with side lengths of order $O(\log H)$. As discussed earlier, the number of permissible units in this lattice region is at most $O((\log H)^{n-1})$.

By multiplying the number of ideals by the number of permissible units, the total number of target irreducible polynomials within the subfield $F$ is bounded by $O(H(\log H)^{n-1})$. Finally, by summing this bound over the finite set of subfields $\mathrm{Sub}_n(K)$, and noting that the cardinality of this set is $O(1)$ relative to $H$, the total number of permissible irreducible polynomials is bounded as follows:
\begin{equation*}
    \#\mathrm{irr}_n(H) \ll \sum_{F \in \mathrm{Sub}_n(K)} H(\log H)^{n-1} \ll H(\log H)^{n-1}.
\end{equation*}

Having determined the cardinalities of both components, the upper natural density of the entire target family is given by:
\begin{equation*}
    \overline{d}(M_K) = \limsup_{H \to \infty} \frac{\#\mathrm{red}_n(H) + \#\mathrm{irr}_n(H)}{\#\mathcal{B}_n(H)} \ll \limsup_{H \to \infty} \frac{O(H^{n-1}) + O(H(\log H)^{n-1})}{O(H^n)}.
\end{equation*}
Consequently, the above limit tends to zero. Since the natural density is a non-negative quantity, it follows that the natural density exists and is equal to zero.
\end{proof}

\end{document}
